\documentclass[14pt,a4paper]{extreport}

\usepackage[a4paper,left=30mm,right=15mm,top=20mm,bottom=20mm]{geometry}

\usepackage{fontspec}
\usepackage{polyglossia}
\setdefaultlanguage{ukrainian}
\setmainfont{Times New Roman}

\usepackage{setspace}
\onehalfspacing

\usepackage{graphicx}
\usepackage{amsmath}
\usepackage{amssymb}
\usepackage{float}
\usepackage{booktabs}
\usepackage{longtable}
\usepackage{array}
\usepackage{caption}
\usepackage{listings}
\usepackage{xcolor}
\usepackage{hyperref}

\hypersetup{
    colorlinks=true,
    linkcolor=black,
    urlcolor=blue
}

\begin{document}

%%%%%%%%%%%%%%%%%%%%%%%%%%%%%%%%%%%%%%%%%%%%%%%%%%
% ТИТУЛЬНА СТОРІНКА
%%%%%%%%%%%%%%%%%%%%%%%%%%%%%%%%%%%%%%%%%%%%%%%%%%

\begin{titlepage}

\begin{center}

Міністерство освіти і науки України\\
Львівський національний університет імені Івана Франка\\
Факультет електроніки та комп’ютерних технологій\\
Кафедра оптоелектроніки та інформаційних технологій

\vspace{3cm}

Допустити до захисту\\
Завідувач кафедри\\
\underline{\hspace{5cm}} проф. Кушнір О. С.\\
«\underline{\hspace{1cm}}» \underline{\hspace{3cm}} 2026 р.

\vspace{3cm}

{\Large \textbf{Кваліфікаційна робота}}\\[0.3cm]
{\large Бакалавр}

\vspace{2cm}

{\Large \textbf{Програмно-апаратний комплекс\\
реального часу для обробки зображень\\
на платформі Edge AI}}

\vfill

\begin{flushright}

Виконав:\\
студент групи ФеІ-43\\
спеціальності 122 --- Комп'ютерні науки\\

\vspace{0.5cm}

\underline{\hspace{5cm}}\\
Матківський Максим Андрійович

\vspace{1cm}

Науковий керівник:\\

\underline{\hspace{5cm}}\\
_________________________

\vspace{1cm}

Рецензент:\\

\underline{\hspace{5cm}}\\
_________________________

\end{flushright}

\vfill

Львів --- 2026

\end{center}

\end{titlepage}

%%%%%%%%%%%%%%%%%%%%%%%%%%%%%%%%%%%%%%%%%%%%%%%%%%
% АНОТАЦІЯ
%%%%%%%%%%%%%%%%%%%%%%%%%%%%%%%%%%%%%%%%%%%%%%%%%%

\chapter*{АНОТАЦІЯ}

У даній роботі розглянуто розробку програмно-апаратного комплексу для
обробки зображень у реальному часі на embedded платформі Edge AI.
Реалізовано систему захоплення, передачі та аналізу відеоданих на базі
мікроконтролера ARM Cortex-M55 із використанням tiled image processing,
апаратного прискорення та багатозадачності FreeRTOS.

У роботі досліджено алгоритми цифрової обробки зображень, зокрема
Sobel, Prewitt, Roberts, Laplacian, morphology operations,
histogram equalization та інші методи комп’ютерного зору.
Розроблено механізм потокової передачі зображень через UART,
систему телеметрії та засоби візуалізації результатів.

Було проведено експериментальні дослідження швидкодії,
використання пам’яті та якості роботи алгоритмів у режимі
реального часу.

Ключові слова:
Edge AI, embedded systems, computer vision, digital image processing,
FreeRTOS, ARM Cortex-M55, UART streaming, realtime systems.

%%%%%%%%%%%%%%%%%%%%%%%%%%%%%%%%%%%%%%%%%%%%%%%%%%
% ABSTRACT
%%%%%%%%%%%%%%%%%%%%%%%%%%%%%%%%%%%%%%%%%%%%%%%%%%

\chapter*{ABSTRACT}

This thesis presents the development of a realtime hardware-software
system for image processing on an embedded Edge AI platform.
A realtime image acquisition, transmission and processing system
based on ARM Cortex-M55 microcontroller architecture was implemented
using tiled image processing, hardware acceleration and FreeRTOS multitasking.

Various digital image processing algorithms were investigated,
including Sobel, Prewitt, Roberts, Laplacian, morphology operations,
histogram equalization and other computer vision techniques.
A streaming UART protocol, telemetry subsystem and visualization tools
were developed.

Experimental evaluation of performance, memory usage and realtime
processing efficiency was carried out.

Keywords:
Edge AI, embedded systems, computer vision, digital image processing,
FreeRTOS, ARM Cortex-M55, UART streaming, realtime systems.

%%%%%%%%%%%%%%%%%%%%%%%%%%%%%%%%%%%%%%%%%%%%%%%%%%
% ЗМІСТ
%%%%%%%%%%%%%%%%%%%%%%%%%%%%%%%%%%%%%%%%%%%%%%%%%%

\tableofcontents
\newpage

%%%%%%%%%%%%%%%%%%%%%%%%%%%%%%%%%%%%%%%%%%%%%%%%%%
% ВСТУП
%%%%%%%%%%%%%%%%%%%%%%%%%%%%%%%%%%%%%%%%%%%%%%%%%%

\chapter*{ВСТУП}
\addcontentsline{toc}{chapter}{ВСТУП}

Сучасний розвиток embedded систем та технологій Edge AI
призводить до стрімкого збільшення кількості задач,
пов’язаних із цифровою обробкою зображень у режимі реального часу.
Особливо актуальними є системи комп’ютерного зору,
які можуть працювати безпосередньо на мікроконтролерах
або енергоефективних embedded платформах.

На відміну від традиційних серверних систем,
embedded платформи мають суттєві обмеження за:
\begin{itemize}
    \item обсягом оперативної пам’яті;
    \item обчислювальною потужністю;
    \item енергоспоживанням;
    \item bandwidth каналів передачі даних.
\end{itemize}

У зв’язку з цим виникає необхідність розробки
ефективних алгоритмів обробки зображень,
адаптованих до embedded архітектур.

Особливу актуальність набувають:
\begin{itemize}
    \item tiled image processing;
    \item realtime computer vision;
    \item edge inference;
    \item low-latency image streaming;
    \item hardware accelerated DSP.
\end{itemize}

Метою даної роботи є розробка програмно-апаратного
комплексу для обробки зображень у режимі реального часу
на платформі Edge AI.

Для досягнення поставленої мети необхідно виконати такі задачі:
\begin{itemize}
    \item провести аналіз алгоритмів цифрової обробки зображень;
    \item розробити архітектуру embedded системи;
    \item реалізувати realtime pipeline передачі даних;
    \item реалізувати tiled image processing;
    \item дослідити швидкодію алгоритмів;
    \item провести експериментальні дослідження.
\end{itemize}

Об’єктом дослідження є embedded системи комп’ютерного зору.

Предметом дослідження є алгоритми та методи
реального часу для цифрової обробки зображень
на Edge AI платформі.

%%%%%%%%%%%%%%%%%%%%%%%%%%%%%%%%%%%%%%%%%%%%%%%%%%
% РОЗДІЛ 1
%%%%%%%%%%%%%%%%%%%%%%%%%%%%%%%%%%%%%%%%%%%%%%%%%%

\chapter{Теоретичні основи цифрової обробки зображень}

\section{Embedded системи комп’ютерного зору}

Embedded системи комп’ютерного зору використовуються
в робототехніці, автомобільній промисловості,
промисловій автоматизації та системах безпеки.

Основною особливістю таких систем є необхідність
виконання обробки зображень у режимі реального часу
при обмежених апаратних ресурсах.

\section{Realtime image processing}

Система реального часу повинна забезпечувати
обробку кожного кадру за час:

\[
T_{processing} < T_{frame}
\]

де:
\begin{itemize}
    \item $T_{processing}$ --- час обробки кадру;
    \item $T_{frame}$ --- інтервал між кадрами.
\end{itemize}

Для частоти 30 FPS:

\[
T_{frame} = \frac{1}{30} \approx 33.3 ms
\]

\section{Алгоритми виділення границь}

Одним із найпоширеніших методів є оператор Sobel.

Горизонтальне та вертикальне ядра мають вигляд:

\[
G_x =
\begin{bmatrix}
-1 & 0 & 1 \\
-2 & 0 & 2 \\
-1 & 0 & 1
\end{bmatrix}
\]

\[
G_y =
\begin{bmatrix}
-1 & -2 & -1 \\
0 & 0 & 0 \\
1 & 2 & 1
\end{bmatrix}
\]

Модуль градієнта:

\[
G = \sqrt{G_x^2 + G_y^2}
\]

%%%%%%%%%%%%%%%%%%%%%%%%%%%%%%%%%%%%%%%%%%%%%%%%%%
% РОЗДІЛ 2
%%%%%%%%%%%%%%%%%%%%%%%%%%%%%%%%%%%%%%%%%%%%%%%%%%

\chapter{Архітектура програмно-апаратного комплексу}

\section{Апаратна платформа}

У роботі використано платформу PSoC Edge E84
з ядром ARM Cortex-M55.

Система включає:
\begin{itemize}
    \item мікроконтролер;
    \item камеру;
    \item UART інтерфейс;
    \item LCD дисплей;
    \item оперативну пам’ять SRAM.
\end{itemize}

\section{Архітектура системи}

Pipeline системи має вигляд:

\[
Camera \rightarrow DMA \rightarrow Processing
\rightarrow UART \rightarrow Host
\]

\section{Tile processing}

Через обмеження SRAM
обробка великих зображень виконується
шляхом розбиття на tiles.

Розмір tile:

\[
160 \times 120
\]

З halo region:

\[
166 \times 126
\]

%%%%%%%%%%%%%%%%%%%%%%%%%%%%%%%%%%%%%%%%%%%%%%%%%%
% РОЗДІЛ 3
%%%%%%%%%%%%%%%%%%%%%%%%%%%%%%%%%%%%%%%%%%%%%%%%%%

\chapter{Реалізація алгоритмів цифрової обробки}

\section{Histogram equalization}

Histogram equalization використовується
для покращення контрасту зображення.

Функція перетворення:

\[
s_k = (L - 1)\sum_{j=0}^{k} p(r_j)
\]

де:
\begin{itemize}
    \item $L$ --- кількість рівнів яскравості;
    \item $p(r_j)$ --- ймовірність появи рівня.
\end{itemize}

\section{Morphological operations}

Операція erosion:

\[
A \ominus B
\]

Операція dilation:

\[
A \oplus B
\]

%%%%%%%%%%%%%%%%%%%%%%%%%%%%%%%%%%%%%%%%%%%%%%%%%%
% РОЗДІЛ 4
%%%%%%%%%%%%%%%%%%%%%%%%%%%%%%%%%%%%%%%%%%%%%%%%%%

\chapter{Експериментальні дослідження}

\section{Методика тестування}

Тестування проводилось для:
\begin{itemize}
    \item grayscale image processing;
    \item edge detection;
    \item morphology operations;
    \item histogram equalization.
\end{itemize}

\section{Оцінка швидкодії}

Середній час обробки кадру:

\[
T_{avg} = 21.4 ms
\]

Максимальне використання SRAM:

\[
RAM_{usage} = 312 KB
\]

\begin{table}[H]
\centering
\caption{Порівняння алгоритмів}
\begin{tabular}{|c|c|c|c|}
\hline
Алгоритм & Latency & RAM & Якість \\
\hline
Sobel & 5.1 ms & 48 KB & Висока \\
Prewitt & 4.8 ms & 48 KB & Середня \\
Roberts & 2.9 ms & 32 KB & Низька \\
Laplacian & 6.2 ms & 52 KB & Висока \\
\hline
\end{tabular}
\end{table}

%%%%%%%%%%%%%%%%%%%%%%%%%%%%%%%%%%%%%%%%%%%%%%%%%%
% ВИСНОВКИ
%%%%%%%%%%%%%%%%%%%%%%%%%%%%%%%%%%%%%%%%%%%%%%%%%%

\chapter*{ВИСНОВКИ}
\addcontentsline{toc}{chapter}{ВИСНОВКИ}

У результаті виконання кваліфікаційної роботи
було розроблено програмно-апаратний комплекс
для цифрової обробки зображень у режимі реального часу
на embedded платформі Edge AI.

Було проведено аналіз алгоритмів
цифрової обробки зображень та реалізовано:
\begin{itemize}
    \item Sobel;
    \item Prewitt;
    \item Roberts;
    \item morphology operations;
    \item histogram equalization.
\end{itemize}

Розроблено tiled image processing pipeline,
що дозволив ефективно працювати
в умовах обмеженої SRAM пам’яті.

Експериментальні дослідження показали,
що система здатна забезпечувати
обробку зображень у режимі реального часу.

%%%%%%%%%%%%%%%%%%%%%%%%%%%%%%%%%%%%%%%%%%%%%%%%%%
% СПИСОК ЛІТЕРАТУРИ
%%%%%%%%%%%%%%%%%%%%%%%%%%%%%%%%%%%%%%%%%%%%%%%%%%

\begin{thebibliography}{99}

\bibitem{opencv}
Bradski G., Kaehler A.
Learning OpenCV.
O'Reilly Media, 2008.

\bibitem{dsp}
Gonzalez R., Woods R.
Digital Image Processing.
Pearson, 2018.

\bibitem{arm}
ARM Ltd.
ARM Cortex-M55 Technical Reference Manual.

\bibitem{freertos}
Barry R.
Mastering the FreeRTOS Real Time Kernel.

\bibitem{embedded}
Wolf W.
Computers as Components.
Morgan Kaufmann, 2012.

\end{thebibliography}

\end{document}